\subsection{\texorpdfstring{Disintegration of a measure \(\mu\) relative to a pseudo-image of \(\mu\)}{Disintegration of a measure mu relative to a pseudo-image of mu}}

THEOREM 2. --- Let T and B be two locally compact spaces having countable bases, \(\mu\) a positive measure on T, \(p\) a \(\mu\) -measurable mapping of T into B, and \(\nu\) a pseudo-image measure of \(\mu\) under \(p\) . Then there exists a \(\nu\) -adequate family \(b \mapsto \lambda_b\) ( \(b \in \mathbf{B}\) ) of positive measures on T, having the following properties:

\begin{enumerate}
\def\labelenumi{\alph{enumi})}
\item
  \(\lambda_{b} \neq 0\) for \(b \in p(\mathrm{T})\) ;
\item
  \(\lambda_{b}\) is concentrated on the set \(\overline{p}^{-1}(b)\) for every \(b\in \mathbf{B}\) ; in particular, \(\lambda_{b} = 0\) for \(b\notin p(\mathrm{T})\) ;
\item
  \(\mu = \int \lambda_b d\nu(b)\)
\end{enumerate}

Moreover, if \(\nu' = r \cdot \nu\) is a second pseudo-image measure of \(\mu\) under \(p\) , and if \(b \mapsto \lambda_b'\) is a \(\nu'\) -adequate family of positive measures on T having the properties b) and c) with respect to \(\nu'\) , then \(\lambda_b = r(b)\lambda_b'\) almost everywhere in B (for \(\nu\) or \(\nu'\) ).

There exists a continuous and finite numerical function \(f\) defined on \(T\) , such that \(f(t) > 0\) for every \(t \in T\) and such that \(\mu'' = f \cdot \mu\) is bounded (Ch. V, §5, No.~6, Prop. 11). Let \(\nu'' = p(\mu'')\) , which is equivalent to \(\nu\) , and write \(\nu'' = g \cdot \nu\) , with \(g\) finite and locally \(\nu\) -integrable; one can suppose, moreover, that \(g(b) > 0\) for all \(b \in B\) (Ch. V, §5, No.~6, Prop. 10). Th. 1 of No.~1, applied to \(\mu''\) and \(\nu''\) , shows that there exists a \(\nu''\) -adequate family \(b \mapsto \lambda_b''\) ( \(b \in B\) ) of positive measures on \(T\) , such that: 1) \(\| \lambda_b'' \| = 1\) for \(b \in p(T)\) ; 2) \(\lambda_b''\) is concentrated on \(\overline{p}^{-1}(b)\) for every \(b \in B\) ; 3) \(\mu'' = \int \lambda_b'' d\nu''(b)\) .

For every \(b \in B\) , let us define a positive measure \(\lambda_{b}\) on T by the formula \(\lambda_{b} = (1/f) \cdot (g(b)\lambda_{b}^{\prime\prime})\) . It is clear that the family \(b \mapsto \lambda_{b}\) has the properties a) and b) of the statement. On the other hand, for every function \(h \in \mathcal{K}(T)\) , h/f belongs to \(\mathcal{K}(T)\) , therefore

\[
\int h (t) d \mu (t) = \int \left(h (t) / f (t)\right) d \mu^ {\prime \prime} (t) = \int d \nu^ {\prime \prime} (b) \int \left(h (t) / f (t)\right) d \lambda_ {b} ^ {\prime \prime} (t).
\]

But since the function \(b \mapsto \int (h(t) / f(t)) d\lambda_b''(t)\) is \(\nu''\) -integrable, the function \(b \mapsto g(b) \int (h(t) / f(t)) d\lambda_b''(t)\) is \(\nu\) -integrable (Ch. V, §5, No.~3, Th. 1).

By the definition of \(\lambda_b\) , this function is \(b \mapsto \int h(t) d\lambda_b(t)\) , whence (loc. cit.) \(\int h(t) d\mu(t) = \int d\nu(b) \int h(t) d\lambda_b(t)\) , which proves that \(\mu = \int \lambda_b d\nu(b)\) .

To establish the second part of the statement, we remark that one can suppose that \(r(b) > 0\) for all \(b \in B\) (Ch. V, §5, No.~6, Prop. 10); set \(\lambda_b^{\prime \prime \prime} = f \cdot \left( \left( r(b) / g(b) \right) \lambda_b' \right)\) ; one shows, as above, that for every function \(h \in \mathcal{K}(T)\) , the relation

\[
\int h (t) d \mu (t) = \int d \nu^ {\prime} (b) \int h (t) d \lambda_ {b} ^ {\prime} (t)
\]

implies

\[
\int h (t) d \mu (t) = \int d \nu^ {\prime \prime} (b) \int \left(h (t) / f (t)\right) d \lambda_ {b} ^ {\prime \prime \prime} (t).
\]

Therefore Th. 1 of No.~1, applied to \(\mu''\) and \(\nu''\) , implies that for almost every \(b \in B\) , \(\lambda_b''' = \lambda_b''\) , whence \(\lambda_b = r(b)\lambda_b'\) .

DEFINITION 2. --- Let T and B be two Polish locally compact spaces. Given a positive measure \(\mu\) on T, a \(\mu\) -measurable mapping \(p\) of T into B, and a pseudo-image measure \(\nu\) of \(\mu\) under \(p\) , every \(\nu\) -adequate family \(b \mapsto \lambda_b\) ( \(b \in B\) ) of positive measures on T having the properties \(b\) ) and \(c\) ) of Th. 2 is called a disintegration of \(\mu\) relative to \(\nu\) .

When p is \(\mu\) -proper and \(\nu = p(\mu)\) , the concept of disintegration relative to p thus coincides with the concept of disintegration relative to \(\nu\) . Under the hypotheses of Th. 2, two disintegrations of \(\mu\) relative to the same pseudo-image measure \(\nu\) are equal almost everywhere for \(\nu\) .

Remark. --- Th. 1 of Ch. V, §3, No.~4 shows (taking into account the fact that T and B have countable bases) that for every function f defined on T, with values in \(\overline{R}\) or in a Banach space F and \(\mu\) -integrable, the set of \(b \in B\) such that f is not \(\lambda_{b}\) -integrable is \(\nu\) -negligible, the function \(b \mapsto \int \mathbf{f}(t) d\lambda_{b}(t)\) , defined almost everywhere, is \(\nu\) -integrable, and

\[
\int \mathbf {f} (t) d \mu (t) = \int d \nu (b) \int \mathbf {f} (t) d \lambda_ {b} (t).
\]

An analogous result holds for scalarly \(\mu\) -integrable functions, on applying Prop. 3 of §1, No.~1.

